% Generated draft: requires mathematical and source review.
\documentclass[11pt]{article}
\usepackage[margin=25mm]{geometry}
\usepackage{fontspec}
\setmainfont{Latin Modern Roman}
\usepackage{amsmath,amssymb,mathtools,unicode-math}
\setmathfont{Latin Modern Math}
\usepackage{enumitem,needspace,xurl,hyperref}
\setlength{\parindent}{0pt}
\setlength{\parskip}{5pt}
\setlength{\emergencystretch}{3em}
\newcommand{\N}{\mathbb{N}}
\newcommand{\Z}{\mathbb{Z}}
\newcommand{\Q}{\mathbb{Q}}
\newcommand{\R}{\mathbb{R}}
\newcommand{\C}{\mathbb{C}}
\newcommand{\F}{\mathbb{F}}
\newcommand{\A}{\mathbb{A}}
\newcommand{\PP}{\mathbb{P}}
\newcommand{\E}{\mathbb{E}}
\newcommand{\eps}{\varepsilon}
\newcommand{\dd}{\,\mathrm{d}}
\begin{document}
\section*{UnsolvedMath ID: 30004664}
\textbf{Source problem code:} catalogue entry
\subsection{Definitions and mathematical statement}
The supplied question concerns redundancies in a recursively generated system of inequalities for a string cone associated with a reduced word \(\mathbf i\). The underlying algebraic data, the permitted reduced words, and the recursive system itself are not specified in the record. Consequently, the following makes the meaning of redundancy precise conditionally on those data being supplied; it does not specify a unique string-cone presentation.

For a fixed intended presentation, let \(V_{\mathbf i}\) be its finite-dimensional real coordinate space, let \(A_{\mathbf i}\) be the finite index set of its defining inequalities, and let \(\ell_{\mathbf i,a}\colon V_{\mathbf i}\to\R\), for \(a\in A_{\mathbf i}\), be their linear forms. Write the intended cone as
\[
C_{\mathbf i}=\{x\in V_{\mathbf i}:\ell_{\mathbf i,a}(x)\geq 0\text{ for every }a\in A_{\mathbf i}\}.
\]
This notation describes a homogeneous linear presentation with non-strict inequalities; the exact intended presentation must be recovered from the supplied references. For each \(a\in A_{\mathbf i}\), set
\[
C_{\mathbf i}^{(-a)}=\{x\in V_{\mathbf i}:\ell_{\mathbf i,b}(x)\geq 0\text{ for every }b\in A_{\mathbf i}\setminus\{a\}\}.
\]
The indexed inequality \(a\) is redundant if deleting it leaves the cone unchanged, that is, if \(C_{\mathbf i}^{(-a)}=C_{\mathbf i}\). Define
\[
R_{\mathbf i}=\{a\in A_{\mathbf i}:C_{\mathbf i}^{(-a)}=C_{\mathbf i}\}.
\]
\textbf{Problem, conditional on the missing presentation data.} For every reduced word \(\mathbf i\) in the intended class, determine \(\lvert R_{\mathbf i}\rvert\) and give a classification of the inequalities indexed by \(R_{\mathbf i}\) in terms of \(\mathbf i\) and the recursive construction.

The count refers to individual deletability in the fixed indexed presentation. It is presentation-dependent, and it need not equal the number of inequalities that can be removed simultaneously. The source record does not settle which counting convention or notion of classification is intended.
\subsection{Short English statement}
For each reduced word \(\mathbf i\), determine which inequalities in the intended recursively generated description of its string cone can be deleted without changing the cone, count them, and classify them. The underlying algebraic setting and the precise inequality system are missing, so this is a conditional statement requiring source clarification and human mathematical review.
\subsection{Sources}
\begin{itemize}
\item Local UnsolvedMath registry, \texttt{data/unsolvedmath/problems.json}, permanent ID \(30004664\), source problem code \texttt{catalogue entry}, title: \emph{Redundant Inequalities in String Cones}. The supplied record attributes the problem to Oberwolfach Reports workshop discussions, gives proposed year \(2021\), and names no individual proposer.
\item Supplied report citation: \emph{Algebraic Groups (2022)}, Oberwolfach Reports, pp.\ 1087--1148, DOI \texttt{10.4171/owr/2021/20}. \url{https://doi.org/10.4171/owr/2021/20}. The supplied literature list also labels this reference \emph{Oberwolfach Report 20/2021}; the differing year labels are retained without reconciliation.
\item Supplied literature reference: Koshevoy--Schumann, \emph{redundancy in string-cone inequalities} (descriptive label supplied by the registry). \url{https://arxiv.org/abs/2109.14439}.
\item Literature assessment dated exactly \texttt{2026-08-22}: the supplied curation labels the problem open and states that the general redundancy-classification question remains unresolved, with structural reductions and special cases established. This is supplied dated curation derived from the OWR material, not an independently verified literature assessment. The record explicitly requests independent verification of time-sensitive status and cited claims.
\end{itemize}
\end{document}
